\section{The sharp rotational-radial benchmark}\label{sec:rotational}
We now solve the rotational-radial subclass exactly.  Let $\mathcal R_R$ be the class of solids obtained by revolving a simple rectifiable arc
\[
\gamma(s)=(\varrho(s),z(s)),\qquad 0\le s\le L\le R,
\]
in the half-plane $\varrho\ge0$, parametrized by arclength, with
\[
\varrho(0)=\varrho(L)=0,
\qquad
\varrho(s)>0\quad(0<s<L).
\]
By \cref{lem:radial}, every member of $\mathcal R_R$ is wrappable by $D_R$ via the direct radial map.  Define
\[
W_{\rm rev}^{\rm rad}(R):=
\sup\{\Vol(B):B\in\mathcal R_R\}.
\]
No convexity assumption is imposed on the meridian or on the resulting solid.

The Mylar variational problem has a long literature~\cite{Paulsen1994,MladenovOprea2003}.  The sharp inequality below is classical in substance; we record a short global proof adapted to the wrapping notation because it identifies the exact barrier that our non-axisymmetric construction will cross.

\begin{theorem}[Sharp rotational-radial bound]\label[theorem]{thm:rotational}
Let $B$ be an embedded solid of revolution generated by a simple north--south meridian of length $L$.  Then
\begin{equation}\label{eq:rotbound}
\Vol(B)\le \frac{\pi}{12I^2}L^3,
\qquad
I:=\int_0^1\frac{du}{\sqrt{1-u^4}}.
\end{equation}
Equality is attained by the classical Mylar profile.  Consequently
\begin{equation}\label{eq:Wrev}
W_{\rm rev}^{\rm rad}(R)
=\frac{\pi}{12I^2}R^3
=\frac43\left(\frac{\Gamma(3/4)}{\Gamma(1/4)}\right)^2R^3
=0.152315527\ldots R^3.
\end{equation}
\end{theorem}

\begin{proof}
Write the arclength-parametrized meridian as
\[
\gamma(s)=(\varrho(s),z(s)),\qquad 0\le s\le L,
\]
and set
\[
a:=\max_{0\le s\le L}\varrho(s).
\]
The coordinate functions are Lipschitz and
\[
\varrho'(s)^2+z'(s)^2=1
\]
for almost every $s$.

Let $\Omega$ be the planar region enclosed by the meridian together with the segment of the rotation axis joining its endpoints.  By the shell formula followed by Green's theorem, and using that the axis segment has $\varrho=0$,
\[
\Vol(B)
=2\pi\iint_{\Omega}\varrho\,d\varrho\,dz
=\pi\left|\int_0^L\varrho(s)^2z'(s)\,ds\right|.
\]
Therefore
\begin{equation}\label{eq:volumeabs}
\frac{\Vol(B)}{\pi}
\le\int_0^L\varrho^2|z'|\,ds
=\int_0^L\varrho^2\sqrt{1-\varrho'^2}\,ds.
\end{equation}
For every $0\le\varrho\le a$, Cauchy--Schwarz gives the pointwise inequality
\begin{equation}\label{eq:pointwiseCS}
\varrho^2\sqrt{1-\varrho'^2}
+|\varrho'|\sqrt{a^4-\varrho^4}
\le a^2.
\end{equation}
Integrating \eqref{eq:pointwiseCS} and using \eqref{eq:volumeabs},
\begin{equation}\label{eq:beforeTV}
\frac{\Vol(B)}{\pi}
\le a^2L-
\int_0^L|\varrho'|\sqrt{a^4-\varrho^4}\,ds.
\end{equation}

Convexity is not needed to estimate the last term.  Put
\[
G(u):=\int_0^u\sqrt{a^4-t^4}\,dt.
\]
Then $G$ is increasing, $G(0)=0$, and $G(a)=\int_0^a\sqrt{a^4-t^4}\,dt$.  Since $\varrho$ starts at $0$, reaches $a$, and returns to $0$, the total variation of $G\circ\varrho$ is at least $2G(a)$.  Hence
\begin{align*}
\int_0^L|\varrho'|\sqrt{a^4-\varrho^4}\,ds
&=\operatorname{Var}(G\circ\varrho)\\
&\ge2\int_0^a\sqrt{a^4-u^4}\,du\\
&=2Ka^3,
\end{align*}
where
\[
K:=\int_0^1\sqrt{1-u^4}\,du.
\]
Thus
\begin{equation}\label{eq:cubicbound}
\frac{\Vol(B)}{\pi}\le a^2L-2Ka^3.
\end{equation}
The right-hand side is maximized over $a\ge0$ at
\[
a=\frac{L}{3K},
\]
giving
\[
\Vol(B)\le\frac{\pi}{27K^2}L^3.
\]
Since
\[
K=\frac23 I,
\]
this is exactly \eqref{eq:rotbound}.

For equality, the absolute-value step in \eqref{eq:volumeabs} forces $z'$ to have a constant sign almost everywhere where $\varrho>0$; equality in the total-variation estimate forbids any extra backtracking of $\varrho$; and equality in \eqref{eq:pointwiseCS} gives
\[
|\varrho'|=\sqrt{1-(\varrho/a)^4},
\qquad
|z'|=\frac{\varrho^2}{a^2}.
\]
Thus the equality profile is the classical Mylar meridian.  Integrating from a pole to the equator and back gives
\[
L=2a\int_0^1\frac{du}{\sqrt{1-u^4}}=2aI,
\]
so $a=L/(2I)=L/(3K)$ and equality is indeed attained.  Taking $L=R$ proves \eqref{eq:Wrev}.
\end{proof}

For comparison, the sphere with north--south meridian length $R$ has volume
\[
V_{\rm sph}(R)=\frac{4}{3\pi^2}R^3
=0.1350949115\ldots R^3.
\]
The three quantitative reference values are therefore
\begin{center}
\begin{tabular}{@{}lc@{}}
\toprule
class or construction & volume divided by $R^3$\\
\midrule
sphere & $0.1350949115\ldots$\\
sharp rotational-radial optimum (Mylar) & $0.1523155270\ldots$\\
certified non-axisymmetric construction & $0.1798434942\ldots$\\
\bottomrule
\end{tabular}
\end{center}
Thus allowing arbitrary nonconvexity while retaining the rotational-radial structure still cannot beat Mylar.  The gain in \cref{thm:main} requires leaving that entire class.

